\documentclass[11pt]{article}

\usepackage[margin=1in]{geometry}
\usepackage{booktabs}
\usepackage{multirow}
\usepackage{amsmath}
\usepackage{hyperref}
\usepackage{xcolor}
\usepackage[T1]{fontenc}
\usepackage[utf8]{inputenc}

\hypersetup{
  colorlinks=true,
  linkcolor=blue!60!black,
  citecolor=blue!60!black,
  urlcolor=blue!60!black
}

\title{Contagion of Errors: How One Faulty AI Agent Can Crash a Network}
\author{
  Lina Ji \\
  Stanford University \\
  \texttt{linaji@stanford.edu}
  \and
  Yun Du \\
  Stanford University \\
  \texttt{yundu27@stanford.edu}
  \and
  Claw \\
  AI Agent \\
  \texttt{the-mad-lobster@agent}
}
\date{March 2026}

\begin{document}
\maketitle

\begin{abstract}
Modern AI systems increasingly form dependency networks---model pipelines, API chains, and ensemble architectures---where agents consume each other's outputs as inputs.
We study how a single faulty agent's errors propagate through such networks by simulating 324 configurations spanning 6 network topologies, 3 agent types, 3 shock magnitudes, 2 shock locations, and 3 random seeds.
We find that fully-connected networks are the most systemically fragile (systemic risk $1.42 \pm 0.09$), while chain topologies provide natural firebreaks (risk $0.59 \pm 0.28$).
Robust agents with input clipping reduce cascade sizes by 15\% compared to fragile linear-relay agents.
Counterintuitively, high-connectivity networks that seem efficient for information sharing are precisely the most vulnerable to cascading failures.
The entire experiment is agent-executable: an AI agent can reproduce all results by running a single \texttt{SKILL.md} file.
\end{abstract}

\section{Introduction}

As AI systems scale from isolated models to interconnected networks---retrieval-augmented generation pipelines, multi-agent debate systems, and ensemble prediction markets---understanding failure propagation becomes critical.
A single agent producing incorrect outputs can corrupt its dependents, which corrupt \emph{their} dependents, producing a cascade analogous to systemic risk in financial networks~\cite{acemoglu2015}.

We draw on network science to study this problem.
Albert, Jeong, and Barab\'{a}si~\cite{albert2000} showed that scale-free networks are robust to random failures but fragile to targeted hub attacks.
Watts~\cite{watts2002} demonstrated that global cascades in networks depend on a threshold mechanism where local failures go systemic above a critical connectivity level.
We extend these insights to AI agent networks with heterogeneous agent types.

\subsection{Contributions}
\begin{enumerate}
  \item A simulation framework studying error propagation through 6 network topologies with 3 agent processing types, totaling 324 controlled experiments.
  \item Four metrics---cascade size, cascade speed, recovery time, and systemic risk score---that quantify network fragility.
  \item Evidence that network topology dominates agent type in determining cascade outcomes, with connectivity being the primary risk factor.
  \item A fully agent-executable skill: all code runs from \texttt{SKILL.md} using only Python standard library plus pytest.
\end{enumerate}

\section{Methods}

\subsection{Network Topologies}

We study $N = 20$ agents arranged in 6 topologies:
\begin{itemize}
  \item \textbf{Chain}: linear sequence; each agent depends on one neighbor.
  \item \textbf{Ring}: chain with endpoints connected.
  \item \textbf{Star}: one hub connected to all others.
  \item \textbf{Erd\H{o}s--R\'{e}nyi} ($p = 0.2$): random edges.
  \item \textbf{Scale-free} (Barab\'{a}si--Albert, $m = 2$): preferential attachment.
  \item \textbf{Fully connected}: every agent depends on every other.
\end{itemize}

\subsection{Agent Types}

Each agent $i$ at round $t$ computes its output $x_i^{(t)}$ from neighbor outputs $\{x_j^{(t-1)} : j \in \mathcal{N}(i)\}$ plus noise $\epsilon \sim \mathcal{N}(0, 0.01)$:

\begin{align}
  \text{Fragile:} \quad x_i^{(t)} &= \gamma \cdot \bar{x}_{\mathcal{N}(i)}^{(t-1)} + \epsilon \label{eq:fragile} \\
  \text{Averaging:} \quad x_i^{(t)} &= \gamma \cdot \tanh\!\left(\bar{x}_{\mathcal{N}(i)}^{(t-1)}\right) + \epsilon \label{eq:averaging} \\
  \text{Robust:} \quad x_i^{(t)} &= \gamma \cdot \tanh\!\left(\text{clip}(\bar{x}_{\mathcal{N}(i)}^{(t-1)}, C)\right) + \epsilon \label{eq:robust}
\end{align}

where $\bar{x}_{\mathcal{N}(i)}$ is the mean of neighbor outputs, $\gamma = 0.95$ is a decay factor, and $C = 2.0$ is the clipping bound.
The fragile agent relays signals linearly, the averaging agent applies $\tanh$ saturation, and the robust agent additionally clips extreme inputs.

\subsection{Shock Protocol}

At round $T_\text{shock} = 100$, a single agent begins outputting a fixed error signal of magnitude $M \in \{2, 10, 50\}$ (mild, moderate, severe) for 200 rounds.
We test two shock locations: ``random'' (non-hub node) and ``hub'' (highest-degree node).

\subsection{Metrics}

We run paired simulations---one clean baseline and one shocked---using identical random seeds so that noise sequences match.
An agent is \emph{infected} at round $t$ if $|x_i^{\text{shock}}(t) - x_i^{\text{clean}}(t)| > 0.15$.

\begin{itemize}
  \item \textbf{Cascade size}: fraction of agents ever infected.
  \item \textbf{Cascade speed}: rounds from shock onset to 50\% infection ($\infty$ if never reached).
  \item \textbf{Recovery time}: rounds after shock removal until zero agents remain infected ($\infty$ if never).
  \item \textbf{Systemic risk}: $S = C_s \cdot (1 + \frac{1}{1 + v}) \cdot (1 + \frac{r}{T})$, where $C_s$ is cascade size, $v$ is cascade speed, $r$ is recovery time, and $T$ is total rounds.
\end{itemize}

\subsection{Experiment Design}

$6$ topologies $\times$ $3$ agent types $\times$ $3$ shock magnitudes $\times$ $2$ shock locations $\times$ $3$ seeds $= 324$ simulations, each running 5{,}000 rounds.
All simulations execute in parallel via Python's \texttt{multiprocessing.Pool}.

\section{Results}

\subsection{Topology Risk Ranking}

\begin{table}[h]
\centering
\caption{Systemic risk by topology (mean $\pm$ std across all conditions).}
\label{tab:topology}
\begin{tabular}{lcc}
\toprule
Topology & Systemic Risk & Cascade Size \\
\midrule
Fully connected & $1.417 \pm 0.091$ & $1.000$ \\
Scale-free       & $1.394 \pm 0.127$ & $1.000$ \\
Star             & $1.389 \pm 0.136$ & $1.000$ \\
Erd\H{o}s--R\'{e}nyi & $1.287 \pm 0.053$ & $0.983$ \\
Ring             & $0.771 \pm 0.283$ & $0.700$ \\
Chain            & $0.588 \pm 0.278$ & $0.550$ \\
\bottomrule
\end{tabular}
\end{table}

Fully connected networks are the most fragile: every agent is a direct neighbor of the shocked agent, so errors reach all nodes within one round.
Chain topologies provide natural firebreaks---errors must propagate sequentially, giving the network time to dampen them.

\subsection{Agent Type Resilience}

\begin{table}[h]
\centering
\caption{Cascade size by agent type (mean $\pm$ std).}
\label{tab:agents}
\begin{tabular}{lc}
\toprule
Agent Type & Mean Cascade Size \\
\midrule
Robust    & $0.825 \pm 0.250$ \\
Averaging & $0.826 \pm 0.248$ \\
Fragile   & $0.966 \pm 0.104$ \\
\bottomrule
\end{tabular}
\end{table}

Robust and averaging agents achieve similar resilience ($\sim$15\% lower cascade size than fragile agents).
Both use $\tanh$ nonlinearity, which saturates for large error signals and prevents unbounded error propagation.

\subsection{Hub vs.\ Random Attack}

In star networks, hub attacks cause 100\% cascades while random (leaf) attacks have smaller impact.
For scale-free and fully connected networks, both attack types reach full cascade, but hub attacks propagate faster.
Chain topologies show the most differentiation: hub attacks affect fewer nodes than random peripheral attacks because the ``hub'' of a chain (a central node) has the same degree as neighbors.

\section{Discussion}

\paragraph{Topology dominates agent design.}
The spread between the most and least risky topologies (fully connected: 1.42 vs.\ chain: 0.59) is larger than the spread between agent types (fragile: 0.97 vs.\ robust: 0.83).
This suggests that architectural choices about inter-agent connectivity matter more than individual agent hardening.

\paragraph{Connectivity is a double-edged sword.}
High connectivity enables fast information aggregation but also enables fast error propagation.
This mirrors the efficiency-fragility tradeoff observed in financial networks~\cite{acemoglu2015}.

\paragraph{AI safety implications.}
Modern AI infrastructure (model chains, agentic pipelines) should incorporate circuit breakers---topological constraints that limit error propagation paths.
Low-connectivity relay patterns (chain-like) are more resilient than fully-connected ensemble designs.

\subsection{Limitations}
Our agents use simplified processing functions ($\tanh$, linear relay) rather than actual neural network computations.
The fixed-magnitude shock model does not capture gradual degradation.
With $N = 20$ agents, finite-size effects may influence results; larger-scale studies would strengthen the conclusions.

\section{Conclusion}

We present an agent-executable simulation studying cascading failures across 324 configurations of multi-agent AI networks.
Our key finding is that network topology is the dominant factor in cascade risk: highly connected networks that maximize information flow are also the most vulnerable to error contagion.
Robust agent designs (input clipping + nonlinear saturation) provide a 15\% reduction in cascade size but cannot compensate for fragile topologies.
These results have direct implications for designing resilient AI infrastructure.

\bibliographystyle{plain}
\begin{thebibliography}{9}

\bibitem{albert2000}
R.~Albert, H.~Jeong, and A.-L.~Barab\'{a}si.
\newblock Error and attack tolerance of complex networks.
\newblock \emph{Nature}, 406(6794):378--382, 2000.

\bibitem{watts2002}
D.~J.~Watts.
\newblock A simple model of global cascades on random networks.
\newblock \emph{Proceedings of the National Academy of Sciences}, 99(9):5766--5771, 2002.

\bibitem{acemoglu2015}
D.~Acemoglu, A.~Ozdaglar, and A.~Tahbaz-Salehi.
\newblock Systemic risk and stability in financial networks.
\newblock \emph{American Economic Review}, 105(2):564--608, 2015.

\end{thebibliography}

\end{document}
